\subsection{補階数のテータ積分の交換}\label{sec:exchange-proof}

後の行列式の比較には、二つの異なる群に沿う積分を、同じ群上の関数として
比較する必要がある。ここでは \(W,V_+,V_-\) の階数をそれぞれ
\(n,d_+,d_-\) とし、\(n=d_++d_-\)、\(d_+-d_-=2e\ge0\) とする。
\(W\) の負符号の個数 \(q\) は \(V_+,V_-\) の負符号の個数の和とする。
指標 \(\nu\) の各無限素点の重さを \(\mu_\tau\) とすると、範囲は
\(2q+2\le n\le2\mu_\tau+2q\) である。
\(\nu\) は共役と交わらない分裂素点の空でない有限集合 \(S\) で順分岐し、
その分岐部分を \(\chi\) とする。\(S\cup\bar S\) の因子を除いた
\(\nu\) の底変換の \(L\) 関数を \(L^S(s,\nu)\) と書き、
\begin{equation}\label{eq:exchange-lstar}
 L_*(e)=\left(\prod_{w\in S}\mathrm Nw\right)^{e^2}
             \prod_{j=1}^{e}|L^S(j,\nu)|^2,\qquad
 L_*(-e)=L_*(e)^{-1}
\end{equation}
とする。二つの有限データ \(\Phi,\Psi\) は、可逆な整行列 \(g\) の
逆行列の末尾 \(d_+\) 行 \(R(g)\) と、\(g\) の先頭 \(d_-\) 列 \(C(g)\)
にそれぞれ支持され、
\(\Phi(R(g))=\chi(\det g)\Psi(C(g))\) を満たすと仮定する。

\(V_-\) 側の無限素点データ \(y\) を固定すると、\(V_+\) 側のデータ \(x\) を
法と指標を選ぶ前に取ることができる。両側の積分を
\(I_+(x,\Phi;g)\)、\(I_-(y,\Psi;g)\) と書くと
\begin{equation}\label{eq:exchanged-integrals}
 I_+(x,\Phi;g)=c\, I_-(y,\Psi;g),\qquad
 |c|^2=C L_*(e)>0
\end{equation}
である。ここで \(g\) は \(U(W)\) の点で、\(C>0\) は法・指標に依らない。
Weil 表現の規約では、\(W\) の分裂に使う指標は
\(\gamma=\nu^{-1}\upsilon^{n/2}\) である。
相手方の二つの指標 \(\alpha,\alpha'\) は \(\alpha\alpha'=\gamma\) を満たす。
この条件は階数を含めた分裂指標の積が \(\nu\) の逆底変換になることと同値である。
\leanlinks{\leansrc{New/Determinant/Exchange/TheoremForPartners}{交換恒等式の仮定と規格化}}

\paragraph{双曲平面を付けたところでの比較}
随伴を取って \(V_-\) の符号を反転させ、まず両側を \(W\) から相手方への
核として表す。\(p=d_-\) とし、\(W_1=W\perp H\)、
\(X=W_1\perp H^{p}\) を使う。大きい側の相手は
\(V_-^{\mathrm{opp}}\perp W_1\)、小さい側の相手は \(V_-^{\mathrm{opp}}\) である。
まず \(y\) を \(W_1\) の側へ持ち上げて、零での切片が \(y\) である
Schwartz ベクトル \(y_1\) を取る。
もう一方の出発データとして \((W_1,W_1)\) の Schwartz ベクトル \(v_1\) を置く。
分岐素点では標準データを用いる。これは整補助行列 \(z\) に対し、
行側では \(z\) と変数を重ねた正方行列の、列側では \(z\) と変数の積の
行列式に \(\chi^{-1}\) を適用した関数である。行列式が非単元なら値を零とする。

\(X\) の双曲分解の混合模型で、\(p\) 本の等方列に沿う切片を使う。
無限素点での標準フレームを \(x_0=\binom{I}{0}\)、相手方の形式の行列を \(J\) とする。
標準フレームの近くで、フレーム \(x\) の張る空間をユニタリ変換で
滑らかに動かせる。実際 \(w=(x_0^*Jx_0)^{-1}x_0^*\) と置き、
\[
 A(B)=(Bx_0w-x_0w B^*)J,\qquad
 \mathcal T(B)=\exp(A(B))(B-A(B))
\]
とする。\(A(B)^*J+JA(B)=0\) で、\((B-A(B))x_0=x_0wB^*Jx_0\) である。
\(A(I)=0\) であり、\(\mathcal T\) の \(I\) での微分は恒等写像だから
逆関数定理を適用できる。\(I+(x-x_0)wJ\) に局所逆を適用し、
得た \(B\) から \(k(x)=\exp(A(B))\) とすれば、
\(x\) は \(k(x)x_0\) と同じ空間を張る。

この近傍に非負で \(x_0\) で \(1\) となる滑らかな切断関数 \(\sigma\) を取る。
各無限素点で、\(x\) の張る空間への標準フレームの射影
\(z(x)=x(x^*Jx)^{-1}x^*Jx_0\) を使う。逆行列は切断の近傍だけで必要なので、
滑らかな延長を選んでよい。大きい側の切片は、\(y_1\otimes v_1\) に
\(k(z(x))\) を作用させ、\(\sigma(z(x))\)、移したフレーム \(k(z(x))^{-1}x\)
に対する \(\sigma\)、およびその上の正方ブロックの
共役行列式の \(2\mu_\tau\) 乗を掛けて作る。小さい側の切片では
同じ標準フレームの切断と、正方変数 \(B_\tau\) の
\(|\det B_\tau|^{2(j+\mu_\tau)}\) を掛ける。ただし \(2j=\operatorname{rk}W_1\) で、
反対の偏極での零切片が \(y_1\) になるように規格化する。
実際、その零切片はフレーム方向の Fourier 変換を零で評価したもので、
切断と行列式因子の積の積分は正である。この積分で割ればよい。
コンパクトなフレーム方向と Schwartz の直交方向の積から、
混合模型の Schwartz ベクトルを得る。

大きい側の積分で \(h\in U(V_-^{\mathrm{opp}}\perp W_1)(E_\infty)\) を動かす。
\(z(h^{-1}x_0)\) の切断が非零なところでは、
\(h k(z(h^{-1}x_0))\) は標準フレームの張る部分を保つのでブロック対角である。
その二つのブロックを \(U(V_-^{\mathrm{opp}})\) と \(U(W_1)\) の積分変数に移す。
動かしたフレームは \(\binom{B_\tau}{0}\) となり、テータ汎関数は
二因子に分かれる。\(U(W_1)\) に沿う積分は
\[
 I_0=\int_{[U(W_1)]}
    \theta_{W_1,W_1}\bigl(v_1,\chi^{-1}\circ\det\bigr)(1,h_1)\,dh_1
\]
である。もう一方には小さい側の切片が残る。
二つの切片の行列式因子の比は、\(\det B_\tau\) を無限部分に持つイデール
\(b_\infty\) に対し \(1/(\nu(b_\infty/\bar b_\infty)|b_\infty|_{\mathbb A_E}^{j})\)
であり、核の Levi 作用に含まれる逆の因子と相殺する。
\leanlinks{\leansrc{New/Determinant/Exchange/RealVectors}{滑らかなフレーム移動と切片の周期}}

中心の Fourier 係数を取ると、核の和は指定された Gram 行列を持つフレームの
軌道に限られる。小さい側が表さない無限素点での符号には、
大きい側も上の切断により寄与しない。有限部分が自己双対で符号が合う場合には、
エルミート形式の局所大域原理と Witt の定理で標準の Gram 行列に移せる。
大きい側の標準フレームの固定群は \(U(W_1)\)、小さい側では自明群である。
大きい側を \(U(V_-^{\mathrm{opp}})\times U(W_1)\) の段階で展開し、
上の切片の積分を代入する。こうして二つの Fourier 係数の比は
\(I_0\) と一つの測度の定数、および絶対値 \(1\) の指標値の積となる。
測度の定数は、非負の切断が単位元の近くで正であるため零でない。

自己双対な格子でこの比較が足りる理由は次の通りである。
分岐素点で、列標準データを新たな双曲線の根群で移し、
法 \(\varpi_w\) の行列 \(A\) を走らせる。
差の中心 Fourier 係数に掛かる和は
\[
 \sum_{A}\chi_w^{-1}(\det A)\,
       \psi_w(\operatorname{tr}(TA))
\]
である。原始指標の行列 Gauss 和なので、\(T\) が特異なら零、
\(T\) が可逆なら \(T=I\) の和に行列式指標を掛けたものである。
従って分岐素点で可逆な指数だけを調べればよい。
残る有限素点では、隣接する双曲平面のせん断で中心係数の部分平均を動かし、
強近似でその Gram 行列が整で逆も整な点まで移す。
この運搬では列標準データの左逆のブロックをせん断が固定する。
そこで得られた Gram 行列は整直交基底をもち、
自己双対の場合の係数の比較を適用できる。
中心 Fourier 反転と放物部分群の分解により、標準データ上の二積分が比例する。
\leanlinks{\leansrc{New/Determinant/Exchange/LargeTriple}{Gram 軌道による Fourier 係数の比較}
  \leanlinksep \leansrc{New/Determinant/Exchange/StableSpacesStandardData}{行列 Gauss 和と自己双対格子での判定}}

\paragraph{定数項を取り階数を下げる}
一枚ずつ対応する双曲平面を除く。除いたあとの大きい側の階数と
中心の階数の半分との差を \(u\) とする。大きい側の定数項を混合模型で展開すると、
等方線上の積分は \(s=u+1\) における Tate 積分になる。
標準データの行列式指標は、零ベクトルとその Fourier 変換側の
零ベクトルの項を消す。分岐素点では単元の行列式指標の Gauss 和を取り、
不分岐素点では整な真空に対する幾何級数を取る。
無限素点では等方錐に沿う切片を極大コンパクト群と
イデールの無限部分で積分する。
従って \(u\ge0\) なら、定数項は \(s=u+1\) で規格化した
錐積分をデータとする小さい階数のテータ積分の \(\gamma_u\) 倍である。
\(u<0\) では Poisson 和公式のもう一方を \(s=-u\) で使う。
両場合で、\(a_u>0\) を指標によらない定数として
\begin{equation}\label{eq:exchange-descent}
 |\gamma_u|^2 L_*(u)=a_u L_*(u+1).
\end{equation}
実際、正の引数 \(s\) での Tate 因子の絶対値は、
固定因子を除いて \(|L^S(s,\nu)|\prod_{w\in S}(\mathrm Nw)^{s-1/2}\) であり、
\((u+1)^2-u^2=2(u+1)-1\) と \eqref{eq:exchange-lstar} がこの式を与える。
列側の定数項は、全階数の列が原始的であるため非零の等方列の寄与を持たず、
零での切片の積分に真空の固定係数 \(\kappa\ne0\) を掛けた値になる。
従って一枚除くと交換係数は \(c\kappa/\gamma_u\) になり、
\eqref{eq:exchange-descent} により指標に依らない定数を除いた
\(|c|^2/L_*(u+1)\) は保たれる。
構成の最初に加えた \(p+1\) 枚をすべて取り除いて、列側を目的の切片に戻す。
交換係数の二乗は固定の正の係数を除いて
\(L_*(e)|I_0|^2/L_*(j)\) になる。
\leanlinks{\leansrc{Lib/SiegelWeil/Tower/LargeSideConstantTerm}{定数項の Tate 積分と局所因子}
  \leanlinksep \leansrc{New/Determinant/Exchange/RemovingPlanes}{錐積分と零切片による降下}}

\paragraph{等階数の規格化と非消滅}
\(\chi^{-1}\circ\det\) を有限データとする等階数の積分は一定である。
分岐素点では列の本数が零のため一般線型群がそのデータを固定し、
他の有限素点では標準レベルが固定する。強近似と等方線の根群に沿う
定数項の計算を合わせると、Levi の線の変数への依存が消える。
この等階数積分の絶対値の二乗を \(L_*(\operatorname{rk}W_1/2)\) で割った値を
\(R_{W_1}(v_1)\) とする。

まず正定値の \(J_0\) を階数 \(2k_0\) とし、\(W_1=H\perp J_0\) とする。
固定した指標で \(R_{W_1}\) が全データに対して零と仮定すると、
上の Fourier 係数の展開により
\((H^{2k_0}\perp W_1,J_0^{\mathrm{opp}}\perp W_1)\) の大きい側の
積分はすべての標準データで零となる。
この消滅は一枚ずつ下へ降りる。実際、任意の下段のデータ \(v\) に対し、
等方錐に小さな滑らかな切断を置き、その上の切片を
\(v\) のコンパクト群による移動で作る。上段の Tate 積分には
正の Mellin 因子 \(B(s)\) と下段の積分 \(I(v)\) が分離する。
\(s\ge1\) で得た積の公式は解析接続で \(s>-k_0\) に延びる。
重さ \(\mu_\tau\ge k_0\) により、関数等式の無限素点因子が
この範囲の整数で零を生じない。\(s>1\) での Euler 積と \(s=1\) の非消滅、
および関数等式により、\(s>-k_0\) の整数で \(L\) 因子は非零である。
定数項を取る点 \(u+1>-k_0\) で上段が零なら \(I(v)=0\) となる。
\(2k_0+1\) 枚除くと \((J_0,0)\) に達するが、ここで真空の積分は
零階数の商の正の体積に等しく、矛盾する。

次に交換を \((J_0,J_0,0)\) と \((J_0,0,J_0)\) に順に適用する。
二つの係数 \(c_1,c_2\) の二乗は、指標から独立な因子を除いてそれぞれ
\(L_*(k_0)R_{W_1}(v_1)\)、\(L_*(k_0)^{-1}R_{W_1}(v_1)\) である。
中間の \((J_0,J_0)\) の核を二つの群で積分し、Fubini で順序を変えると、
\(c_1c_2\) は零階数の実データの値に等しく、指標に依らない。
二つの \(L_*(k_0)\) が消えるので \(R_{W_1}(v_1)^2\) も指標に依らない。
先の非消滅から一つの \(v_1\) を選ぶと、\(R_{W_1}(v_1)>0\) は
どの指標でも同じ値になる。
一般の \(W_1\) の符号へは双曲平面を加えて移る。
上段の等階数積分は自身の定数項に等しく、指定した下段の実データを
錐積分で実現できるので、各段で \eqref{eq:exchange-descent} により
\(R_{W_1}\) は指標に依らない正の定数倍だけ変わる。
\leanlinks{\leansrc{New/Determinant/Exchange/ConeImage}{錐積分で指定したデータを実現}
  \leanlinksep \leansrc{New/Determinant/Exchange/VanishingDescends}{定数項による消滅の降下}
  \leanlinksep \leansrc{New/Determinant/Exchange/TheoremForPartners}{等階数の比の非消滅と指標からの独立}}

\paragraph{任意の対応するデータへの拡張}
二積分の差は有限部分の左移動について同変な線型写像である。
核には標準列データ \(P(Y)=\chi^{-1}(\det zY)\) が入っている。
\(z x_0=I\) となる整列 \(x_0\) を取り、列データ \(\Psi\) は
\(\varpi^a\) を法とする合同で値が一定になるように \(a\) を取る。
\(d=x_0 z+\varpi^{-a}(I-x_0 z)\) は \(x_0\) の張る部分を固定し、
\(\ker z\) を拡大する。\(P(dgY)\ne0\) なら
\(Y\equiv g^{-1}x_0(z\,dgY)\pmod{\varpi^a}\) である。
右作用に対する指標 \(\chi^{-1}\circ\det\) から
\(\Psi(g^{-1}x_0)P(dgY)=|P(dgY)|\Psi(Y)\) を得る。
左辺を \(g\in\mathrm{GL}_n(\mathcal O_S)\) で平均すると核の元である。
右辺の \(|P(dgY)|\) の平均は支持上で一定で、\(Y=x_0,g=1\) の近くで正である。
従って \(\Psi\) 自身も核に入る。法 \(\mathfrak m\) のデータを
局所整数環に引き上げてこの議論を適用する。
最後に、小さい側の随伴を戻す際の転置およびノルム \(-1\) の整単元による
変数変換は絶対値 \(1\) の係数だけを生じる。
\(W\) 側の分裂指標の残りも、\(\alpha\alpha'=\gamma\) により両核に同じ
\(\alpha(\det g)^{-1}\) が掛かる。これで \eqref{eq:exchanged-integrals} が従う。
\leanlinks{\leansrc{New/Determinant/Exchange/StableSpacesStandardData}{平行移動の平均による対応データの生成}
  \leanlinksep \leansrc{New/Determinant/Exchange/TheoremForPartners}{随伴と分裂指標の変更}}

\paragraph{持ち上げを二回合成した場合}
\(W=U_1\perp U_2\)、\(V_+=A\)、\(V_-=A'\) とする。
行データ \(x\) の二つの列ブロックを核に置き、\(U_1\) から \(A\) への持ち上げと
\(A\) から \(U_2\) への転置持ち上げを合成して \(\Theta_x\) とする。
列データ \(y\) の二つの行ブロックを使う \(A'\) 経由の合成を \(\Theta'_y\) とする。
比較のため \(U_1\) の形式 \(f_1\) と \(U_2\) の形式 \(f_2\) を固定し、
合成した像を \(f_2\) と対にする。
各辺の二つの分裂指標を \((a_i,b_i)\) と \((a_i',b_i')\) とすると、
積の条件から \(a_1a_2=a_1'a_2'=\gamma\) であり、
\(\nu_0=b_1/b_2=b_2'/b_1'\) となる。
\((g_1,g_2)\) における交換恒等式に
\(f_1(g_1)\overline{f_2(g_2)}\nu_0(\det g_2)\) を掛ける。
この最後の指標は、両経路で二つの核を積に分ける際の指標の差を打ち消す。

\(A'\) 側は、二つの辺の Fock 多項式を最初から積にしておく。
点データも二つの行ブロックに分かれるので、Fubini と持ち上げ・転置の関係から
その積分は \(\langle f_2,\Theta'_y f_1\rangle\) になる。
\(A\) 側の実ベクトルは交換で選ばれたものである。まず多項式を
二つの列ブロックの多項式の積の和に書く。各積について、辺の実型の
再生公式でそれぞれの持ち上げを所定の辺の核の実群平均に直す。
対にする積分でこの二つの平均を合わせると、実データにだけ依存する
一つの係数が残る。両変数の形式は辺の実型のもとで急減少するため、
積分は Schwartz ベクトルのモーメントのノルムについて連続である。
従って実データが一般でも積分は \(\kappa_0\langle f_2,\Theta_x f_1\rangle\) である。
これは、すべての有限データと形式について同じ比例関係が多項式上で成り立ち、
その比例関係が閉じていることによる。
両方を有限データの重みで足して任意の \(f_2\) と対にすると
\begin{equation}\label{eq:exchange-square}
 \sum_y\Psi(y)\Theta'_y=\frac{\kappa_0}{c}\sum_x\Phi(x)\Theta_x .
\end{equation}
ここでは \(c\ne0\) によって割った。\(\kappa_0\) は捻りに依らず、
この段階の式は \(\kappa_0=0\) の場合も含む。
\leanlinks{\leansrc{New/Determinant/Exchange/OneSquare}{二つの転置を含む経路の交換}}
